% ===================================================================
%  CUADRO N.º 6 — La Casa por dentro. — Los Muebles. — La Comida.
%                 El Alumbrado.
%  Traduction 2026 d'apres texto/io, controlee sur texto/fr et
%  texto/en. Consigne : voir texto/es/10-cuadro-01.tex.
%
%  Deux notes, toutes deux marquees « (*) », et deux appels : celui de
%  « babo » au bloc c1-12-1, celui de « lambrequino » au bloc c2-01-1.
%  L'ordre les departage, comme dans les autres colonnes.
%
%  La femme de chambre porte le (16) que gravuri/korekti.json declare
%  pour ce bloc : l'ido imprime (6), le francais (16), et la planche
%  porte bien un 16. Voir texto/en/15-table-06.tex.
% ===================================================================

%%K t06-apar-1 apar
\VUsaut{6.0mm}
\VUtitre{15.0pt}{60}{CUADRO N.º 6}
\VUsaut{1.4mm}
\VUtitre{11.4pt}{0}{La Casa por dentro, los Muebles, la Comida,}
\VUsaut{0.4mm}
\VUtitre{11.4pt}{0}{el Alumbrado.}
\VUsaut{2.2mm}

%%K t06-apar-2 apar
La casa está terminada. El tío de Juan se ha instalado en su nueva
vivienda, cuya distribución ya conocemos. Veamos ahora cómo la ha
amueblado. Visitaremos primero:

%%K t06-tit-1 sub
\VUpk{10.2pt}{40}{El Dormitorio.}
\VUsaut{3.0mm}

%%K t06-c1-01-1 p
1. --- Llegamos a mala hora, porque la criada está ocupada limpiando la
\VUgras{habitación}.

%%K t06-c1-02-1 p
2. --- Sobre el \VUgras{lavabo} \textsuperscript{(1)} están esparcidos los distintos
objetos con los que uno se lava, se peina, en una palabra, se arregla.
Son la \VUgras{palangana} \textsuperscript{(2)} y el \VUgras{jarro} \textsuperscript{(3)}, el
\VUgras{cepillo del pelo} \textsuperscript{(4)}, el \VUgras{peine} \textsuperscript{(5)}, el
\VUgras{jabón} \textsuperscript{(6)} y la \VUgras{esponja} \textsuperscript{(7)}; por último,
un \VUgras{frasco} \textsuperscript{(8)} de elixir dentífrico.

%%K t06-c1-03-1 p
3. --- En el suelo, delante del lavabo está el \VUgras{cubo}
\textsuperscript{(9)} para recoger las aguas sucias.

%%K t06-c1-04-1 p
4. --- En la \VUgras{antesala} \textsuperscript{(10)} vemos algunas \VUgras{perchas}
\textsuperscript{(13)} y dos \VUgras{paraguas} \textsuperscript{(11)} en el
\VUgras{paragüero} \textsuperscript{(12)}.

%%K t06-c1-05-1 p
5. --- Al otro lado de la puerta está el \VUgras{armario de luna}
\textsuperscript{(14)}, que guarda la \VUgras{ropa blanca} \textsuperscript{(15)}.

%%K t06-c1-06-1 p
6. --- Aprovechemos que la \VUgras{doncella} \textsuperscript{(16)} hace la
\VUgras{cama} \textsuperscript{(17)} para estudiar sus distintas partes. El
\VUgras{armazón de la cama} \textsuperscript{(20)} lleva un buen \VUgras{somier}
\textsuperscript{(21)} sobre el que Justina va a poner enseguida un
\VUgras{colchón} \textsuperscript{(22)} de lana. Luego cogerá de las sillas las dos
\VUgras{sábanas} \textsuperscript{(23)} y la \VUgras{manta} \textsuperscript{(24)}. Después
colocará a la cabecera de la cama el \VUgras{travesero} \textsuperscript{(25)} y la
\VUgras{almohada} \textsuperscript{(26)}, que están con el \VUgras{edredón}
\textsuperscript{(27)} sobre la \VUgras{alfombrilla de cama} \textsuperscript{(32)}. Un
plumero para sacudir el polvo de las \VUgras{cortinas} \textsuperscript{(19)} y del
\VUgras{dosel} \textsuperscript{(18)}, y ya está hecha una parte del trabajo.

%%K t06-c1-07-1 p
7. --- Luego tendrá que ordenar un poco la \VUgras{mesilla de noche}
\textsuperscript{(28)}, dar cuerda al \VUgras{despertador} \textsuperscript{(29)}, preparar
la \VUgras{lamparilla} \textsuperscript{(30)} y llevarse la \VUgras{vela}
\textsuperscript{(31)} para limpiar la palmatoria.

%%K t06-tit-2 sub
\VUsaut{5.0mm}
\VUpk{10.2pt}{40}{El costurero y los accesorios de la chimenea.}
\VUsaut{3.0mm}

%%K t06-c1-08-1 p
8. --- El \VUgras{costurero} \textsuperscript{(34)} que está junto al \VUgras{taburete}
\textsuperscript{(33)} también necesita mucho que lo ordenen. Habrá que doblar el
\VUgras{bordado} \textsuperscript{(35)}, guardar en la \VUgras{mesa de labor}
\textsuperscript{(38)} el \VUgras{dedal}, el \VUgras{hilo} \textsuperscript{(36)}, las
\VUgras{agujas} \textsuperscript{(37)} y las \VUgras{tijeras} \textsuperscript{(39)}, y cerrar
con la \VUgras{llave} \textsuperscript{(40)} la tapa de la mesa.

%%K t06-c1-09-1 p
9. --- En el \VUgras{velador} \textsuperscript{(41)} del centro, Justina renovará el
agua de la \VUgras{garrafa} \textsuperscript{(42)}. Pondrá \VUgras{azúcar} en el
\VUgras{azucarero} \textsuperscript{(43)}, limpiará las \VUgras{pinzas del azúcar}
\textsuperscript{(44)} y se llevará el \VUgras{cepillo de la ropa} \textsuperscript{(46)}.

%%K t06-c1-10-1 p
10. --- El lado derecho de la habitación le dará menos trabajo, porque la
\VUgras{tumbona} \textsuperscript{(47)} y su \VUgras{cojín} \textsuperscript{(48)} están en
su sitio y, como no hace frío, nada se ha movido entre los accesorios de
la \VUgras{chimenea} \textsuperscript{(49)}. \VUgras{Pala} \textsuperscript{(50)},
\VUgras{tenazas} \textsuperscript{(51)}, \VUgras{morillos} \textsuperscript{(55)} y
\VUgras{guardacenizas} \textsuperscript{(56)} están limpios; la \VUgras{pantalla de
chimenea} \textsuperscript{(54)} no se ha movido y el \VUgras{cajón de la leña}
\textsuperscript{(52)} está bien provisto de \VUgras{leña} \textsuperscript{(53)}.

%%K t06-noto-1 noto
\VUnotes{143.2mm}{%
(*) Babo = niño pequeño; inglés \textit{baby}, francés \textit{bébé},
italiano \textit{bambino}.%
}

%%K t06-noto-2 noto
\VUnotes{143.2mm}{%
(*) Lambrequino = francés \textit{lambrequin}, español
\textit{lambrequines}, portugués \textit{lambrequins}, e incluso alemán e
inglés \textit{lambrequins}; es decir, alemán, inglés, francés, portugués
y español por igual.%
}

%%K t06-c1-11-1 p
11. --- Tendrá que quitar además el polvo del \VUgras{reloj de sobremesa}
\textsuperscript{(57)}, de los \VUgras{candelabros} \textsuperscript{(58)} y de los
\VUgras{vaciabolsillos} \textsuperscript{(59)}, y luego limpiar el \VUgras{espejo}
\textsuperscript{(60)}.

%%K t06-c1-12-1 p
12. --- En cuanto a la \VUgras{cuna} \textsuperscript{(61)} del niño (del babo (*)),
ya está lista.

%%K t06-tit-3 sub
\VUsaut{5.0mm}
\VUpk{10.2pt}{40}{El Salón.}
\VUsaut{3.0mm}

%%K t06-c2-01-1 p
1. --- Pasemos ahora al salón. Es una \VUgras{habitación} muy bonita. La
chimenea, que está en el rincón de la derecha, va adornada con una fina
\VUgras{estatuilla} \textsuperscript{(1)} y dos hermosos \VUgras{jarrones}
\textsuperscript{(3)}, y drapeada con un \VUgras{lambrequín} \textsuperscript{(2)} de
terciopelo (*).

%%K t06-c2-02-1 p
2. --- Para evitar que las chispas del hogar prendan fuego a la alfombra,
se ha puesto delante un \VUgras{guardafuegos} \textsuperscript{(4)} de cobre.

%%K t06-c2-03-1 p
3. --- Al lado, un elegante \VUgras{escritorio} \textsuperscript{(5)} de señora está
abierto. Allí es donde la \VUgras{señora de la casa} \textsuperscript{(23)} escribe
sus cartas.

%%K t06-c2-04-1 p
4. --- La ventana va guarnecida con \VUgras{visillos} \textsuperscript{(6)} recogidos
en \VUgras{alzapaños} \textsuperscript{(8)} colgados de las \VUgras{escarpias}
\textsuperscript{(7)}, y con grandes cortinas de tela rematadas por una
\VUgras{cenefa} \textsuperscript{(9)}.

%%K t06-c2-05-1 p
5. --- En el hueco de la ventana, un \VUgras{macetero} \textsuperscript{(11)} contiene
una \VUgras{planta} \textsuperscript{(10)} verde, una palmera magnífica cuyas hojas
llegan casi hasta la \VUgras{lámpara de petróleo} \textsuperscript{(12)}.

%%K t06-c2-06-1 p
6. --- La \VUgras{pantalla} \textsuperscript{(13)} de la lámpara la ha hecho María,
la encantadora \VUgras{joven} \textsuperscript{(14)} que está al \VUgras{piano}
\textsuperscript{(15)}. Sentada muy derecha en el \VUgras{taburete} \textsuperscript{(16)},
estudia una pieza. Es muy hábil y muy cuidadosa, como lo prueba su
\VUgras{mueble de la música} \textsuperscript{(17)}, que está perfectamente
ordenado.

%%K t06-tit-4 sub
\VUsaut{5.0mm}
\VUpk{10.2pt}{40}{El Travieso.}
\VUsaut{3.0mm}

%%K t06-c2-07-1 p
7. --- Su hermano Pablo busca un \VUgras{libro} \textsuperscript{(19)} en la
\VUgras{estantería} \textsuperscript{(18)}. Es un \VUgras{muchacho} \textsuperscript{(20)}
inquieto e insoportable. Al colgarse de la estantería para alcanzar el
libro que está demasiado alto, se arriesga a tirar el \VUgras{busto}
\textsuperscript{(21)} que hay encima.

%%K t06-c2-08-1 p
8. --- Aprovecha, para hacer esa tontería, que su madre está ocupada
charlando con una amiga, una \VUgras{señora} \textsuperscript{(22)} que ha venido a
pasar unos días con ella. La madre está sentada en el \VUgras{sofá}
\textsuperscript{(24)}, al lado del \VUgras{sillón} \textsuperscript{(25)}, y su amiga está
en la \VUgras{silla} \textsuperscript{(27)}, de la que solo vemos el
\VUgras{respaldo} \textsuperscript{(26)}.

%%K t06-c2-09-1 p
9. --- El gran \VUgras{marco} \textsuperscript{(30)} colgado de la pared contiene el
\VUgras{retrato} \textsuperscript{(29)} de una parienta. En el \VUgras{álbum}
\textsuperscript{(32)} que está sobre la mesa se reúnen las \VUgras{fotografías}
\textsuperscript{(33)} de los demás miembros de la familia. Los niños lo han
hojeado hace un momento, porque las \VUgras{sillas} \textsuperscript{(31)} siguen
agrupadas alrededor de la mesa.

%%K t06-c2-10-1 p
10. --- El \VUgras{jarrón chino} \textsuperscript{(34)} de porcelana que está junto
al \VUgras{abrecartas} \textsuperscript{(35)} se lo regalaron a María por su santo,
y el \VUgras{biombo} \textsuperscript{(36)} que adorna el rincón de la habitación lo
trajo de China su tío.

%%K t06-c2-11-1 p
11. --- Alegrado por los hermosos \VUgras{cuadros} \textsuperscript{(37)} colgados
del \VUgras{tabique} \textsuperscript{(42)}, iluminado por la \VUgras{lámpara de
techo} \textsuperscript{(38)}, cuyas \VUgras{bombillas} \textsuperscript{(39)} dan por la
noche una luz tan alegre, este salón parece muy agradable. La mullida
\VUgras{alfombra} \textsuperscript{(40)} extendida sobre el entarimado y el
\VUgras{timbre eléctrico} \textsuperscript{(41)}, tan práctico, acaban de hacerlo
muy cómodo.

%%K t06-tit-5 sub
\VUsaut{3.6mm}
\VUpk{10.2pt}{40}{El Comedor.}
\VUsaut{3.0mm}

%%K t06-c3-01-1 p
1. --- Ha sonado la hora de la cena. Toda la familia, salvo María, está
reunida alrededor de la mesa. El comedor está agradablemente caldeado por
una excelente \VUgras{estufa} \textsuperscript{(1)} de loza, con un \VUgras{tubo}
\textsuperscript{(2)} delgado.

%%K t06-c3-02-1 p
2. --- Esta sala no tiene muchos muebles. A lo largo de la pared cuelga,
encima del \VUgras{escabel} \textsuperscript{(4)}, una \VUgras{fuentecilla}
\textsuperscript{(3)} decorativa. Muy cerca está el \VUgras{aparador} \textsuperscript{(5)},
sobre el que se dejan los platos fríos, los postres y los distintos
utensilios de mesa que de momento no hacen falta.

%%K t06-c3-03-1 p
3. --- Así vemos, en la balda de arriba, un \VUgras{pollo asado}
\textsuperscript{(7)}, un \VUgras{pescado} \textsuperscript{(8)} y una \VUgras{copa}
\textsuperscript{(10)} con \VUgras{fruta} \textsuperscript{(9)}. Debajo están el
\VUgras{queso} \textsuperscript{(11)}, el \VUgras{pan} \textsuperscript{(12)} y las
\VUgras{vinagreras} \textsuperscript{(13)}, que llevan todo lo necesario para aliñar
la ensalada: el aceite, el vinagre, la \VUgras{sal} \textsuperscript{(14)} y la
\VUgras{pimienta} \textsuperscript{(15)}. Por último, en la balda de abajo hay un
\VUgras{pastel} \textsuperscript{(17)} junto al \VUgras{mostacero} \textsuperscript{(16)}.

%%K t06-c3-04-1 p
4. --- El reloj del comedor, que se llama \VUgras{reloj de pared}
\textsuperscript{(20)}, tiene una forma muy particular. Va encajado en un marco de
madera que contiene además un \VUgras{barómetro} \textsuperscript{(19)} y un
\VUgras{termómetro} \textsuperscript{(18)}.

%%K t06-tit-6 sub
\VUsaut{4.6mm}
\VUpk{10.2pt}{40}{El Servicio de la cena.}
\VUsaut{3.0mm}

%%K t06-c3-05-1 p
5. --- Todas las paredes de la sala van revestidas de \VUgras{paneles de
madera} \textsuperscript{(21)} de roble encerado. Una \VUgras{portera} \textsuperscript{(22)}
de tela cierra bastante bien la puerta que da a la galería. Allí irá la
familia a tomar el café después de cenar. Las \VUgras{tazas} \textsuperscript{(24)} y
los \VUgras{platillos} \textsuperscript{(25)} ya están preparados en el
\VUgras{aparador de la vajilla} \textsuperscript{(23)}.

%%K t06-c3-06-1 p
6. --- Encima de la mesa hay fijada al techo una \VUgras{lámpara colgante}
\textsuperscript{(26)}, cuya luz muy viva alumbra por la noche todo el comedor.

%%K t06-c3-07-1 p
7. --- Ocupémonos ahora de la \VUgras{mesa del comedor} \textsuperscript{(27)} misma.
Cuando entró la familia, la mesa estaba puesta. Sobre el \VUgras{mantel}
\textsuperscript{(28)} se habían colocado, en el sitio de cada comensal, un
\VUgras{plato} \textsuperscript{(33)}, una \VUgras{servilleta} \textsuperscript{(29)}, una
\VUgras{cuchara} \textsuperscript{(34)}, un \VUgras{tenedor} \textsuperscript{(35)}, un
\VUgras{cuchillo} \textsuperscript{(36)} y un \VUgras{vaso} \textsuperscript{(32)}. Había
también \VUgras{botellas} \textsuperscript{(30)} para el vino y una \VUgras{garrafa}
\textsuperscript{(31)} para el agua.

%%K t06-c3-08-1 p
8. --- El \VUgras{criado} \textsuperscript{(39)} trajo primero la \VUgras{sopera}
\textsuperscript{(37)}, en la que todavía humea un poco de sopa. Ahora sirve el
primer \VUgras{plato} \textsuperscript{(38)}, un plato de carne. Luego presentará los
que ya hemos nombrado y, por fin, después del \VUgras{postre}, aprovechará
que sus señores habrán pasado a la galería para quitar la mesa y poner en
orden el comedor.

%%K t06-c3-08-2 p
\textit{Nota.} --- \textit{Las palabras corrientes relativas a la comida
se dan aquí de un modo muy sumario. La lista se completará en los dos
cuadros «El Hotel» (n.º 13) y «El Mercado» (n.º 15).}

%%K t06-tit-7 sub
\VUsaut{4.6mm}
\VUpk{10.2pt}{40}{La Cocina.}
\VUsaut{3.0mm}

%%K t06-c4-01-1 p
1. --- En la cocina, en la que echamos una ojeada por la puerta
entreabierta, encontramos un desorden pintoresco. Delante del
\VUgras{hogar} \textsuperscript{(1)}, donde chisporrotea el \VUgras{fuego}
\textsuperscript{(3)}, está montado el \VUgras{asador giratorio} \textsuperscript{(5)}. Un
hermoso \VUgras{asado} \textsuperscript{(7)} está en el \VUgras{espetón}
\textsuperscript{(6)} y acaba de hacerse.

%%K t06-c4-02-1 p
2. --- El \VUgras{fuelle} \textsuperscript{(8)} y la \VUgras{pala del carbón}
\textsuperscript{(9)}, que se acaban de usar, se han quedado en el suelo, lo mismo
que los \VUgras{leños} \textsuperscript{(10)}.

%%K t06-c4-03-1 p
3. --- La repisa de la chimenea está ordenada: el \VUgras{farol}
\textsuperscript{(11)}, el \VUgras{cerillero} \textsuperscript{(13)}, con las
\VUgras{cerillas} \textsuperscript{(12)} dentro, el \VUgras{candelero}
\textsuperscript{(14)} y la \VUgras{palmatoria} \textsuperscript{(15)} con su \VUgras{vela}
\textsuperscript{(16)} están en su sitio. Pero el gran \VUgras{cubo del carbón}
\textsuperscript{(18)}, lleno del \VUgras{carbón de piedra} \textsuperscript{(17)} que la
\VUgras{cocinera} \textsuperscript{(23)} acaba de ir a buscar a la bodega, se ha
quedado en medio de la cocina.

%%K t06-c4-04-1 p
4. --- La verdad es que la pobre mujer tiene que darse prisa para que la
comida esté lista a tiempo. Ahora está junto al \VUgras{fogón}
\textsuperscript{(19)}, dando vueltas a una \VUgras{salsa} \textsuperscript{(24)} que no
debe dejar que se queme.

%%K t06-c4-05-1 p
5. --- En el \VUgras{horno} \textsuperscript{(20)} se cuece un pastel que ha
preparado para la merienda de los niños. Encima del fogón cuelgan el
\VUgras{cucharón} \textsuperscript{(21)} y la \VUgras{espumadera} \textsuperscript{(22)}.

%%K t06-tit-8 sub
\VUsaut{4.0mm}
\VUpk{10.2pt}{40}{La Batería de cocina.}
\VUsaut{3.0mm}

%%K t06-c4-06-1 p
6. --- La \VUgras{batería de cocina} \textsuperscript{(25)}, colgada al fondo de la
sala, está muy limpia. Las hermosas \VUgras{cacerolas} \textsuperscript{(26)} de
cobre, el \VUgras{lechero} \textsuperscript{(27)}, el \VUgras{colador} \textsuperscript{(28)} y
el \VUgras{rallador} \textsuperscript{(29)} se han limpiado con esmero.

%%K t06-c4-07-1 p
7. --- La \VUgras{salera} \textsuperscript{(30)} está puesta sobre el pequeño
aparador, al lado de la \VUgras{tetera} \textsuperscript{(31)} y de la
\VUgras{damajuana del aceite} \textsuperscript{(32)}. El \VUgras{cajón}
\textsuperscript{(33)} guarda seguramente los cubiertos y los cuchillos de cocina.

%%K t06-c4-08-1 p
8. --- La \VUgras{escoba} \textsuperscript{(34)} y el \VUgras{plumero} \textsuperscript{(35)}
están en un rincón. La cocinera acaba de usar el \VUgras{tajo}
\textsuperscript{(36)}; lo ha dejado junto a la ventana.

%%K t06-c4-09-1 p
9. --- Una \VUgras{toalla de manos} \textsuperscript{(37)} cuelga de la pared, al
lado del \VUgras{embudo} \textsuperscript{(38)}. En esta parte de la cocina se
guardan la \VUgras{parrilla} \textsuperscript{(39)}, la \VUgras{cuchilla de picar}
\textsuperscript{(40)} y la \VUgras{tabla de picar} \textsuperscript{(41)}. Luego, encima
de la cuchilla, en un \VUgras{estante} \textsuperscript{(42)}, hay \VUgras{pucheros}
\textsuperscript{(43)}, la \VUgras{olla} \textsuperscript{(44)} con su \VUgras{tapadera}
\textsuperscript{(45)} y la \VUgras{marmita} \textsuperscript{(46)}. La \VUgras{sartén}
\textsuperscript{(47)} cuelga al lado.

%%K t06-c4-10-1 p
10. --- Solo el \VUgras{grifo} \textsuperscript{(48)} de cobre pone un brillo en este
rincón. El \VUgras{fregadero} \textsuperscript{(49)}, sobre el que está el gran
\VUgras{cántaro} \textsuperscript{(50)}, debe de ser muy cómodo con su armarito, en
el que se guardan el \VUgras{tonel del vinagre} \textsuperscript{(51)} con su
\VUgras{espita} \textsuperscript{(52)}, la \VUgras{caja de la basura} \textsuperscript{(53)} y
la \VUgras{vajilla sucia} \textsuperscript{(54)}.

%%K t06-tit-9 sub
\VUsaut{4.0mm}
\VUpk{10.2pt}{40}{Otras cosas que se guardan en la cocina.}
\VUsaut{3.0mm}

%%K t06-c4-11-1 p
11. --- El gran \VUgras{barreño} \textsuperscript{(55)} en el que se lavan las
\VUgras{verduras} \textsuperscript{(58)} y el \VUgras{caldero} \textsuperscript{(56)} de
hierro colado puesto sobre el \VUgras{embaldosado} \textsuperscript{(57)} también
deben de tener su sitio en ese armario.

%%K t06-c4-12-1 p
12. --- A la cocinera no le faltan hornillos, porque hay además, en la
esquina de la mesa, un \VUgras{hornillo de gas} \textsuperscript{(59)} en el que se
calienta agua en una cacerola pequeña. El \VUgras{vapor} \textsuperscript{(60)} que
sale de ella nos indica que debe de estar hirviendo. Seguramente esa agua
servirá para el café, porque en el otro extremo de la mesa están la
\VUgras{cafetera} \textsuperscript{(64)} y el \VUgras{molinillo de café} \textsuperscript{(63)}.

%%K t06-c4-13-1 p
13. --- El hornillo está unido al \VUgras{mechero de gas} \textsuperscript{(62)} por
un largo \VUgras{tubo de goma} \textsuperscript{(61)}.

%%K t06-c4-14-1 p
14. --- La \VUgras{asistenta} \textsuperscript{(66)}, que viene a ayudar a la
cocinera, prepara las verduras para la cena. Cuando estén peladas y
lavadas, las pondrá en la \VUgras{cazuela} \textsuperscript{(65)} hasta la hora de
cocerlas.

%%K t06-tit-10 sub
\VUsaut{4.0mm}
{\centering\fontsize{10.2pt}{10.2pt}\selectfont\textls[40]{\scshape El Cuarto de Baño.} \textit{(La sala del baño.)}\par}
\VUsaut{3.0mm}

%%K t06-c5-01-1 p
1. --- Solo nos queda una indiscreción por cometer para conocer las
principales habitaciones de la casa: ver la instalación del cuarto de
baño. No será largo, porque no es grande, y sabremos pronto que la
\VUgras{bañera} \textsuperscript{(1)} tiene un buen \VUgras{calentador}
\textsuperscript{(2)}, que la \VUgras{jabonera} \textsuperscript{(3)} está al alcance de la
mano del que se baña y que el \VUgras{toallero} \textsuperscript{(5)} permite secar
muy fácilmente las \VUgras{toallas} \textsuperscript{(4)}.

%%K t06-c5-02-1 p
2. --- Esta sala tiene los tabiques cubiertos de \VUgras{azulejos}
\textsuperscript{(6)}, lo que ha permitido instalar una \VUgras{ducha}
\textsuperscript{(7)} sin temer que el agua que salpica al usarla estropee las
paredes. Una pequeña \VUgras{palangana} \textsuperscript{(8)} de cinc, que sirve
para los baños de pies, completa el mobiliario.
\VUsaut{6.0mm}
\VUfilet{22mm}
